% EGA V, English translation by P. Blass, 1987 (cote 24): archive cote 24, site page 18 (PDF page 19).
% typescript. Draft compared: none.
% AI-transcribed by Claude Opus 5.5 (claude-opus-5-5, adaptive thinking, effort high), Message Batches API, batch pilot (msgbatch_0189Nhf1xB615YWnztubtLfc), 2026-11-12; not yet checked by eye.
\scanpage{24/18}
\mspage{18}

not be immersed it is necessary and sufficient that $Z_i$ should be
such [PB incomprehensible].* If that is so then the radicial
multiplicity of $F$ at $Z_i$ at $k$ is equal to that of $G_\eta$ at $Z_i$ relative
to $k(\eta)$.

The first assertion is contained in the last assertion of
2.7. For the second, since $Y_\eta \rightarrow X$ is flat therefore the functor $F$
$\rightarrow G_\eta$ is exact, we are reduced by restriction to a neighborhood of
the general point of $Z_i$ and by a devissage (unscrewing) to the case
where $F = O_{Z_i}$ and we may assume $Z_i = X$. Also we could start by
assuming that $X$ is separable over $k$ is reduced to the case of $k$
algebraically closed [PB: incomprehensible].* Then the assertion
is contained in 2.15. (iii). Then we conclude as usual:

\uline{Corollary 2.17}

Let $Z$ be the set of points of $X$ where $F$ is not separable over
$k$ (reference). Then $G_\eta$ is separable over $k(\eta)$ if and only if $Z$ is
finite. In particular if $F$ is separable over $k$ then $G_\eta$ is
separable over $k(\eta)$.

\uline{Remark 2.18}

The key result 2.12 (and its consequences 2.13 and 2.17)
become false if we abandon the assumption that $f$ is an immersion
as we see for example by taking for $X$ a purely inseparable
covering of $\underline{P}$. However if $k$ is of characteristic zero 2.12 and
2.17 are valid without assuming that $f$ is an immersion.
